\documentclass[11pt,a4paper]{article}

% ── Packages ─────────────────────────────────────────────────────────────────
\usepackage[T1]{fontenc}
\usepackage[utf8]{inputenc}
\usepackage{mathptmx}
\usepackage{amsmath,amssymb,amsthm}
\usepackage{amsfonts}
\AtBeginDocument{\let\hbar\relax
  \DeclareMathSymbol{\hbar}{\mathord}{AMSb}{"7E}}
\usepackage{geometry}
\usepackage{microtype}
\usepackage{hyperref}
\usepackage{booktabs}
\usepackage{enumitem}
\usepackage{setspace}
\usepackage{titlesec}
\usepackage{abstract}
\usepackage{natbib}

\geometry{top=1.1in, bottom=1.1in, left=1.15in, right=1.15in}

\titleformat{\section}{\normalfont\bfseries\large}{\thesection}{1em}{}
\titleformat{\subsection}{\normalfont\bfseries\normalsize}{\thesubsection}{1em}{}
\titleformat{\subsubsection}{\normalfont\bfseries\itshape\normalsize}{\thesubsubsection}{1em}{}

\newtheorem{theorem}{Theorem}[section]
\newtheorem{proposition}[theorem]{Proposition}
\newtheorem{conjecture}[theorem]{Conjecture}
\newtheorem{definition}[theorem]{Definition}
\newtheorem{remark}[theorem]{Remark}
\newtheorem{postulate}[theorem]{Postulate}

\newcommand{\Pact}{\mathcal{P}_{\mathrm{act}}}
\newcommand{\Tact}{\hat{T}_{\mathrm{act}}}
\newcommand{\ARA}{A_{\mathrm{RA}}}
\newcommand{\SBDG}{S_{\mathrm{BDG}}}

\usepackage{xcolor}
\hypersetup{
  colorlinks=true,
  linkcolor=blue,
  citecolor=blue,
  urlcolor=blue
}

\begin{document}

% ── Title page ───────────────────────────────────────────────────────────────
\begin{titlepage}
\centering
\vspace*{1.5cm}

{\LARGE\bfseries Relational Actualism and the Standard Model\\[0.4em]
\Large\itshape Particles, Forces, and the Topology of the Causal Graph\\[0.4em]
\normalsize\upshape (RASM) --- Prospectus\par}

\vspace{1.2cm}

{\large Joshua F.\ Sandeman\\[0.3em]
\normalsize Independent Researcher, Salem, Oregon, USA\\[0.2em]
\href{mailto:sansuikyo@gmail.com}{sansuikyo@gmail.com}}

\vspace{0.8cm}
{\small\textit{Preprint:}~\href{https://zenodo.org/communities/relational-actualism}%
{zenodo.org/communities/relational-actualism}\par}
\vspace{1.5cm}

\begin{abstract}
Relational Actualism (RA) grounds physics in a single ontological commitment:
that irreversible on-shell actualization events, writing permanent directed
edges into a growing causal Directed Acyclic Graph (DAG), are the primitive
physical reality. The companion papers in this suite have shown that quantum
mechanics is exact at the discrete level, that special-relativistic kinematics
follow from integer counting ($c = l_P/t_P$, $E = \gamma mc^2$), and that
general relativity is the unique consistent macroscopic description of the DAG,
inevitable by Lovelock's theorem given the Lean-verified conservation and
covariance results. This paper extends the program to the Standard Model.
We propose that the Benincasa-Dowker-Glaser (BDG) local topology of a 4D
causal graph has exactly four independent integer-valued degrees of freedom
($N_1, N_2, N_3, N_4$), one of which is already used by gravity.
The remaining \emph{three} give exactly the three non-gravitational gauge
interactions: electromagnetism ($U(1)$), the weak force ($SU(2)$), and the
strong force ($SU(3)$). From this single structural observation we derive:
electric charge quantization in units of $e/3$ (three spatial dimensions plus
the Lean-verified LLC); photon masslessness (zero net BDG debt per cycle);
topological colour confinement and exact baryon number conservation (LLC on
$N_2$ spatial winding); maximal parity violation as a theorem (DAG acyclicity
plus retarded BDG action); all three force ranges from BDG cost and LLC balance;
Regge trajectories; exactly three generations and no fourth ($N_5$ absent in
4D BDG); the Koide lepton mass formula from $SU(3)_{\mathrm{gen}}$ symmetry,
with scale $m_0 \approx m_p/3$ to $0.35\%$; neutrino oscillations from
actualization frequency differences; grand unification as the Planck-scale limit
of $SU(3)_{\mathrm{colour}} \times SU(3)_{\mathrm{gen}}$ giving
$\sin^2\theta_W = 3/8$; and the gauge coupling hierarchy from BDG coefficient
magnitudes. Specific predictions: lightest neutrino mass $m_1 \approx 0.37~\mathrm{meV}$,
$\sum m_\nu \approx 59~\mathrm{meV}$ (testable with CMB-S4/Euclid), and no proton
decay (exact baryon conservation). Supersymmetry is structurally absent.
Three open derivation targets, all requiring the BDG-filter MCMC, are precisely
scoped; when closed, they derive all fermion masses and coupling constants from
the five integers $(1,-1,9,-16,8)$.
Additional consequences: $\theta_{\mathrm{QCD}} = 0$ exactly
(strong CP dissolved; no axion needed); the black hole information
paradox dissolved via Causal Severance partition; and the baryon
asymmetry as a Causal Severance initial condition.
\end{abstract}

\medskip\noindent
\textbf{Keywords:} Relational Actualism \textperiodcentered{}
Standard Model \textperiodcentered{} Causal DAG \textperiodcentered{}
Benincasa-Dowker Action \textperiodcentered{} Topological Particles
\textperiodcentered{} Gauge Symmetry \textperiodcentered{}
Colour Confinement \textperiodcentered{} Three Generations
\textperiodcentered{} Supersymmetry \textperiodcentered{}
Hierarchy Problem \textperiodcentered{} Poisson-CSG
\textperiodcentered{} Causal Bandwidth

\end{titlepage}

% ── 1. Introduction ──────────────────────────────────────────────────────────
\section{Introduction: One Ontology, All of Physics}
\label{sec:intro}

The Standard Model of particle physics is the most precisely tested theory in
science. It describes three of the four known fundamental forces via gauge
fields with symmetry group $U(1) \times SU(2) \times SU(3)$, and organises
the matter content of the universe into three generations of quarks and leptons.
Yet it offers no explanation for why exactly these symmetries, why exactly three
generations, why the specific masses and coupling constants, and why the forces
have the ranges they do. These are inputs, not outputs.

The Relational Actualism suite has so far derived general relativity from first
principles (RAGC~\citep{Sandeman2026RAGC}), explained flat galactic rotation
curves and the Hubble tension without dark matter (RADM~\citep{Sandeman2026RADM}),
and derived quantum error correction thresholds and the speed of light from the
causal graph structure (RAQI~\citep{Sandeman2026RAQI}). In each case, the
explanatory power came from staying entirely within the native ontology: one
growing causal DAG, one actualization criterion, one generating mechanism.

This paper extends the same programme to the Standard Model. The central claim
is simple: the Benincasa-Dowker-Glaser (BDG) local topology of a 4-dimensional
causal graph has exactly four independent integer-valued degrees of freedom.
One linear combination of these is already used by gravity. The remaining
three are the three non-gravitational forces. This is not a postulate; it is
a counting argument from the structure of the BDG action in 4D.

Throughout, we maintain the RA suite's epistemic discipline: established results
are clearly distinguished from conjectures, and open derivation targets are
stated with collaborator-level precision. This paper is a prospectus for a
research programme, not a completed derivation. What it establishes is that the
programme is well-defined, that RA is uniquely positioned to execute it, and
that several structural predictions --- including the absence of supersymmetric
partners --- are already a consequence of the framework.

% ── 2. The Native RA Account of Gravity Reviewed ─────────────────────────────
\section{The Native RA Account of Gravity: What We Already Have}
\label{sec:gravity_review}

Before extending to the Standard Model, it is worth stating precisely what the
RA account of gravity already achieves, because the extension will use the same
mechanism in the same way.

\paragraph{Matter and spacetime as two aspects of the same graph.}
In general relativity, Einstein's equation $G_{\mu\nu} = 8\pi G\,T_{\mu\nu}$
equates the geometry of a smooth manifold with the stress-energy of quantum
fields. Wheeler summarised this as ``matter tells spacetime how to curve;
spacetime tells matter how to move,'' but offered no account of how two
ontologically distinct categories communicate.

In RA, the dualism dissolves: both sides of the equation are aspects of the
same causal DAG. Dense regions of the graph --- those with high local
actualization density $\lambda(x) \gg \bar\lambda$ --- are what we call matter.
Sparse transverse edges connecting these dense regions are what we call
spacetime. There is no matter \emph{inside} spacetime; there is only the graph.
Einstein's equation is the thermodynamic equation of state balancing the
internal complexity of the dense regions against the outward connectivity of
the transverse edges.

\paragraph{The BDG sign criterion and curvature.}
The Benincasa-Dowker-Glaser action~\citep{Benincasa2010} provides the unique
discrete, retarded, Lorentz-invariant action functional for a causal set.
At vertex $v$ in 4D:
\begin{equation}
  \SBDG^{(4)}(v)
  = \frac{1}{l_P^2}\bigl(1 - N_1(v) + 9N_2(v) - 16N_3(v) + 8N_4(v)\bigr),
  \label{eq:BDG_4d}
\end{equation}
where $N_k(v)$ counts the $k$-element inclusive causal intervals in the
immediate past of $v$. The RA actualization criterion $\Delta S(\rho\|\sigma_0) > 0$
has a natural discrete counterpart: vertex $v$ actualizes if and only if
$\SBDG(v) > 0$. Dense matter regions produce $\langle \SBDG \rangle > 0$;
the vacuum has $\langle \SBDG \rangle = 0$.

\paragraph{GR by Lovelock.}
Because the RA-CSG satisfies conservation (the LLC, Lean-verified), covariance
(causal invariance, Lean-verified), and locality (definitional in the CSG
algorithm), Lovelock's theorem~\citep{Lovelock1971} guarantees that the unique
consistent macroscopic field equations are Einstein's. General relativity is
not derived in a limit; it is logically inevitable. The interesting physics
is where RA departs from GR: sparse regions ($\mu \sim 1$), Planck-scale
corrections, Causal Severance boundaries.

% ── 3. The Bandwidth Derivation of Special Relativity ────────────────────────
\section{Causal Bandwidth, the Speed of Light, and Relativistic Kinematics}
\label{sec:bandwidth}

Section 4.4 of RAGC~\citep{Sandeman2026RAGC} establishes the following results
from RA primitives alone. We collect them here because the bandwidth argument
is the bridge between the kinematic and topological structures of RASM.

\paragraph{The speed of light as graph scale ratio.}
The causal graph has two fundamental scales: the Planck length $l_P$ and the
Planck time $t_P$. Their ratio is:
\begin{equation}
  c = \frac{l_P}{t_P}.
  \label{eq:c_native}
\end{equation}
This is not a postulate; it is the maximum rate at which a causal signal can
propagate through the graph. A massless pattern ($m = 0$) has rest-frame
actualization frequency $f_0 = mc^2/h = 0$: no internal bandwidth is consumed
by mass-maintaining actualization, so all bandwidth goes to spatial propagation
at exactly $c$.

\paragraph{Relativistic energy from discrete counting.}
From the on-shell condition and the Lorentz transformation of the actualization
frequency four-vector:
\begin{equation}
  E = \gamma mc^2.
  \label{eq:rel_energy}
\end{equation}
This follows from $f_0 = mc^2/h$, $f_{\mathrm{coord}} = \gamma f_0$, and
$E = hf_{\mathrm{coord}}$ --- no additional postulates. As $v \to c$,
$\gamma \to \infty$: a massive particle would need to consume the entire local
causal bandwidth for spatial propagation while simultaneously maintaining its
on-shell identity, which is impossible at finite energy. This is the RA
account of why massive particles cannot reach $c$.

\paragraph{Proper time as integer count.}
Proper time $\tau = \mathcal{N}(\gamma) \cdot t_P$ is exact at every scale,
including the Planck scale. There is no continuum limit required; the integer
count is the fundamental quantity and the GR proper time integral is its
macroscopic approximation.

% ── 4. The BDG Degrees of Freedom: Four Dimensions, Four Numbers ──────────────
\section{The BDG Degrees of Freedom: Why 4D Gives Exactly Three Forces}
\label{sec:four_dof}

The BDG action in equation~\eqref{eq:BDG_4d} counts four quantities: $N_1, N_2,
N_3, N_4$. These are not independent choices of what to count; they are the
complete set of causal interval sizes that exist in a locally 4-dimensional
graph. The key structural observation of this paper is:

\begin{proposition}[Four degrees of freedom, one used by gravity]
\label{prop:dof}
The BDG local topology of a 4-dimensional causal graph has exactly four
independent integer-valued degrees of freedom $(N_1, N_2, N_3, N_4)$.
One specific linear combination --- $\SBDG^{(4)}$ --- is used by gravity.
The remaining three independent combinations are available as conserved
topological charges of persistent causal patterns.
\end{proposition}

\paragraph{Dimensional dependence.}
In $d$ spacetime dimensions (even $d$), the BDG action counts causal intervals
of sizes $1, 2, \ldots, n_d$ where $n_d = d/2 + 2$~\citep{Benincasa2010}:
\begin{equation}
  d = 2:\quad n_2 = 3 \to N_1, N_2, N_3
  \quad\Rightarrow\quad \text{1 BDG combination, 2 remaining.}
\end{equation}
\begin{equation}
  d = 4:\quad n_4 = 4 \to N_1, N_2, N_3, N_4
  \quad\Rightarrow\quad \text{1 BDG combination, 3 remaining.}
\end{equation}

The 2D case is instructive: with only one non-gravitational combination,
there can be at most one non-gravitational gauge interaction, consistent
with the fact that 1+1D physics supports no propagating gauge degrees of freedom.

The 4D case gives exactly three remaining combinations. This is why
$U(1) \times SU(2) \times SU(3)$ has three factors: not by postulate, but
because the BDG topology of 4D spacetime has exactly three non-gravitational
degrees of freedom. The number of forces is fixed by the spacetime dimension.

\begin{remark}
The spacetime dimension itself is fixed to 4 by the Benincasa-Dowker conjecture
(the BDG action recovers the Einstein-Hilbert action in the continuum limit
only in the physically observed dimension) and by Lovelock's theorem (Einstein
equations require 4D for the standard form). RA therefore predicts $d = 4$ as
the unique self-consistent dimension, and consequently exactly three
non-gravitational forces.
\end{remark}

% ── 5. Particles as Persistent Self-Similar BDG Patterns ─────────────────────
\section{Particles as Persistent Self-Similar BDG Patterns}
\label{sec:particles}

In a growing causal DAG, a particle cannot be a fixed object --- there are no
fixed objects; only the graph grows. A particle must be a \emph{persistent
pattern}: a local causal topology that recreates itself as new vertices are
added under BDG dynamics.

\begin{definition}[Stable particle pattern]
A stable particle pattern is a persistent self-similar local causal topology
in the BDG-filtered Poisson-CSG such that every vertex $v$ in the pattern
satisfies $\SBDG(v) > 0$.
\end{definition}

The BDG stability criterion is the RA actualization filter applied to the
pattern itself: the pattern continuously actualizes, vertex by vertex, so that
it maintains causal presence in the growing graph. An unstable pattern is one
where at some vertex $v$, $\SBDG(v) \leq 0$: the pattern ceases to contribute
to the local curvature and dissipates into the background vacuum. This is
particle decay.

\paragraph{Mass as actualization frequency.}
From the bandwidth result, a pattern of rest mass $m$ actualizes at rest-frame
frequency $f_0 = mc^2/h$. A heavier particle is a more complex BDG pattern
requiring more frequent on-shell actualization events per unit proper time to
maintain its structural stability. Mass is not an intrinsic property of a
label; it is the actualization cost of maintaining a specific causal topology.

\paragraph{Spin from periodicity.}
The periodicity of a self-similar pattern under BDG growth determines its spin:
\begin{itemize}[noitemsep]
  \item \textbf{Bosons (spin 1):} the pattern recreates its local topology
  in one step of graph growth. Integer spin is 1-periodicity.
  \item \textbf{Fermions (spin 1/2):} the pattern alternates between two
  complementary topological states as the graph grows, requiring two steps
  to fully recreate itself. Half-integer spin is 2-periodicity.
\end{itemize}
This discrete account of spin does not rely on the rotation group of a
continuum background; it is a property of the pattern's self-similarity
period in the growing graph.

\paragraph{Anti-particles.}
An anti-particle is a pattern with the same BDG stability and periodicity
as its partner, but with the signs of all three non-gravitational topological
charges reversed. The particle and anti-particle patterns are
mirror-images in the three-dimensional charge space $(Q_1, Q_2, Q_3)$.

\begin{proposition}[Exact baryon number conservation]
\label{prop:baryon}
Baryon number is exactly conserved in RA. A baryon is a pattern carrying
net $N_2$ spatial winding number. The LLC applied to all three spatial
components of $N_2$ forbids any process that changes the net $N_2$ winding
number without a corresponding anti-baryon creation or annihilation.
Since the LLC is Lean-verified (RAGC~\citep{Sandeman2026RAGC}), baryon
number conservation in RA is machine-checked, not merely approximate.
This is a stronger result than the Standard Model, where baryon number
is an accidental symmetry that grand unified theories can violate.
\emph{Testable prediction:} proton decay is forbidden in RA.
\end{proposition}

% ── 6. The Three Forces from the Three Non-Gravitational BDG Combinations ─────
\section{The Three Forces from the Three Non-Gravitational BDG Combinations}
\label{sec:forces}

We identify the three non-gravitational conserved charges of a persistent
BDG pattern with the three gauge charges of the Standard Model. The
identification is structural: each charge is the conserved version, under
the LLC, of one of the three non-gravitational BDG topological invariants.

\subsection{Electromagnetism: Charge Quantization and Photon Masslessness}

The first non-gravitational BDG combination is the signed $N_1$ count: the
net number of immediate causal predecessors, with sign determined by spatial
direction. This is conserved by the LLC at every vertex.

\paragraph{Charge quantization from three spatial dimensions.}
In a 3+1D causal graph, a vertex can receive signed edges from at most three
independent spatial directions. The net signed $N_1$ charge is therefore
constrained to integer values in $\{-3,\,-2,\,-1,\,0,\,+1,\,+2,\,+3\}$.
Identifying the unit charge with $e/3$ (one spatial edge per dimension):
\begin{equation}
  Q_{N_1} \in \{-e,\,-2e/3,\,-e/3,\,0,\,+e/3,\,+2e/3,\,+e\}.
  \label{eq:charge_quantization}
\end{equation}
This matches the complete set of electric charges in the Standard Model:
up quark $+2e/3$ ($Q_{N_1} = +2$), down quark $-e/3$ ($Q_{N_1} = -1$),
electron $-e$ ($Q_{N_1} = -3$), neutrino $0$ ($Q_{N_1} = 0$),
$W^\pm$ ($Q_{N_1} = \pm 3$), and neutral particles ($Q_{N_1} = 0$).
Electric charge is quantized in units of $e/3$ because the spatial
dimensionality of the causal graph is exactly 3.

\paragraph{Photon masslessness from zero net BDG charge.}
The photon carries $Q_{N_1} = Q_{N_2} = Q_{N_3/N_4} = 0$. A pattern with all
non-gravitational BDG charges zero has zero net BDG action per propagation cycle:
\begin{equation}
  \Delta \SBDG^{\text{photon}}\big|_{\text{per cycle}} = 0.
  \label{eq:photon_massless}
\end{equation}
Zero net actualization cost per cycle gives $f_0 = 0$ and therefore $m = 0$.
The photon is massless because its pattern carries no BDG debt of any kind;
all causal bandwidth is allocated to spatial propagation at $c = l_P/t_P$.

\subsection{The Strong Force: Confinement from LLC Spatial Balance}

The second non-gravitational BDG combination measures the medium-range
causal branching of the pattern: how does the pattern's causal ancestry
spread through the two-step neighbourhood $N_2$? In a graph with spatial
structure, this branching has three independent components corresponding to
the three spatial dimensions.

The $N_2$ branching charge has three independent spatial components
$(Q_r, Q_g, Q_b)$, one per spatial dimension. The LLC requires at every vertex:
\begin{equation}
  \Delta Q_r = 0, \qquad \Delta Q_g = 0, \qquad \Delta Q_b = 0.
  \label{eq:colour_llc}
\end{equation}

\paragraph{Confinement from LLC imbalance.}
A single quark pattern carrying colour charge $(Q_r, 0, 0)$ accumulates a
permanent ledger imbalance in the $r$-direction as the pattern propagates.
The causal closure condition
\begin{equation}
  \lim_{t\to\infty}\bigl[Q_r\text{ accumulated}\bigr] = 0
  \label{eq:causal_closure}
\end{equation}
is violated for a free quark. An isolated colour-charged pattern cannot persist
as a stable causal subgraph. This is RA-native confinement: not an external
mechanism but a direct consequence of the Lean-verified LLC.

The only configurations satisfying equation~\eqref{eq:colour_llc} simultaneously
in all three spatial directions are baryons ($r + g + b = 0$) and mesons
($r + \bar{r} = 0$). The $SU(3)$ gauge group is the continuous symmetry of
the LLC constraint $Q_r + Q_g + Q_b = \mathrm{const}$ on $N_2$ spatial charges:
the permutation group $S_3$ of three objects extends continuously to $SU(3)$.

\paragraph{Gluons and asymptotic freedom.}
Gluons are the patterns that mediate exchanges of colour charge. They carry
colour themselves --- their patterns have non-zero $N_2$ branching invariant
--- which is why they interact with each other (unlike photons). At high
energies, the Poisson-CSG parameter $\mu \gg 1$ and the graph is densely
connected: the colour branching modes are so uniformly distributed that they
decouple from any specific pattern, giving asymptotic freedom. At low energies
($\mu \sim 1$), the branching becomes sparse and colour charges clump,
giving the running of $\alpha_s$.

\subsection{The Weak Force: $SU(2)$, Maximal Parity Violation, and the Arrow of Time}

The third non-gravitational BDG combination is the $N_3/N_4$ handedness
charge, encoding how the long-range causal structure of a pattern winds
around the temporal axis as it propagates forward.

\paragraph{Parity violation from DAG acyclicity: a derivation.}
The BDG action is \emph{retarded}: $\SBDG(v)$ counts causal intervals
exclusively in $\mathrm{past}(v)$. The $N_3$ count enumerates three-step
causal chains $a \to b \to c \to v$; the $N_4$ count enumerates four-step
chains $a \to b \to c \to d \to v$. These longer-range intervals encode how
the pattern's causal ancestry spirals around the temporal axis as the pattern
propagates.

For a pattern propagating forward in time, the handedness of its $N_3/N_4$
winding is defined relative to the temporal direction: \emph{left-handed} if
the winding accumulates in the same direction as temporal flow;
\emph{right-handed} if opposite. To convert a left-handed pattern to
right-handed --- to flip the sign of its $N_3/N_4$ winding --- would require
some causal edges to point from future to past. The DAG is acyclic: this is
impossible.

\begin{proposition}[Maximal parity violation from acyclicity]
\label{prop:chirality}
In a DAG with the BDG actualization filter $\SBDG(v) > 0$, the $N_3/N_4$
handedness charge can only accumulate in the direction consistent with forward
temporal propagation. The weak interaction, which changes the $N_3/N_4$
charge of a pattern, can therefore only couple to patterns whose winding is
consistent with forward time --- the left-handed sector. Parity violation is
exactly maximal: not a statistical preference but a topological impossibility
for the right-handed coupling, following directly from DAG acyclicity.
\end{proposition}

In the Standard Model, the left-handedness of the weak interaction is a
postulate with no explanation. In RA it is a theorem: the weak force is
left-handed because time has a direction, the DAG encodes that direction
in its acyclicity, and the retarded BDG action can only read causal structure
from the past.

\paragraph{The arrow of time and the weak force are the same fact.}
The RA account makes precise a long-suspected connection: CP violation in the
weak sector and the arrow of time share a common origin. In RA both are
consequences of DAG acyclicity. The arrow of time is encoded in the
irreversibility of actualization~\citep{Sandeman2026RAGC}; maximal parity
violation is encoded in the retarded BDG action applied to
forward-propagating patterns. They are two faces of the same topological fact.

Weak isospin is the conserved $N_3/N_4$ handedness charge. The $SU(2)$
gauge symmetry acts on doublets of left-handed patterns; right-handed patterns
are $SU(2)$ singlets. \emph{Remaining derivation target:} prove formally that
$N_3/N_4$ winding reversal in a forward-propagating DAG pattern requires a
backward-pointing edge --- a precise combinatorial statement about causal
interval structure.

\paragraph{W and Z bosons as massive exchange patterns.}
The $W$ and $Z$ bosons carry net $N_3/N_4$ handedness charge and therefore
require continuous actualization to persist. Their actualization cost is their
mass, and their short range is the finite proper-time lifetime following from
the non-zero BDG debt per propagation step (Section~\ref{sec:ranges},
equation~\eqref{eq:weak_range}).

\begin{remark}[The Higgs mechanism in RA]
The Higgs pattern mediates the topological operation of adding an $N_3/N_4$
handedness increment to an otherwise neutral pattern, allowing a left-handed
fermion to acquire the $N_3/N_4$ structure that gives it mass. The Higgs
field is the fluctuation mode of the BDG handedness combination around its
vacuum value. This is a conjecture requiring formal development.
\end{remark}

% ── 7. Three Generations from Three BDG Excitation Levels ────────────────────
\section{Three Generations from Three BDG Excitation Levels}
\label{sec:generations}

The BDG action in 4D counts four interval sizes: $N_1, N_2, N_3, N_4$.
The minimum-complexity stable pattern uses only $N_1$ and $N_2$ structure:
this is the \emph{first generation}. More complex patterns incorporating
$N_3$ structure form the \emph{second generation}; those incorporating $N_4$
structure form the \emph{third generation}.

\begin{conjecture}[Three generations from BDG levels]
The three generations of the Standard Model correspond to the three levels
of BDG excitation available in 4D:
\begin{align}
  \text{Generation 1} &: \text{dominant }N_2\text{ internal structure}
  \quad\to\quad e, u, d, \nu_e\\
  \text{Generation 2} &: \text{dominant }N_3\text{ internal structure}
  \quad\to\quad \mu, c, s, \nu_\mu\\
  \text{Generation 3} &: \text{dominant }N_4\text{ internal structure}
  \quad\to\quad \tau, t, b, \nu_\tau
\end{align}
\end{conjecture}

\paragraph{The mass hierarchy.}
Because $N_3$ and $N_4$ structure requires more causal vertices to maintain
(a larger actualization footprint), higher-generation patterns have
correspondingly higher actualization frequencies --- higher masses. The BDG
coefficients $(1, -1, 9, -16, 8)$ determine the relative actualization cost
of each level. The mass ratios of the three generations should therefore
be derivable from these coefficients: a precise calculation connecting the
integers of the BDG action to the observed mass spectrum
(approximately $1 : 207 : 3477$ for charged leptons).

\paragraph{No fourth generation.}
Because the 4D BDG action counts exactly $N_1, N_2, N_3, N_4$ --- four interval
sizes, no more --- there are exactly three excitation levels above the basic
$N_1$ structure. There is no $N_5$ in the 4D BDG action. This is a structural
prediction: \emph{no fourth generation of matter exists}.

This prediction is already supported by the LEP measurement of the number of
light neutrino species ($N_\nu = 2.9840 \pm 0.0082$) and by the absence of
any fourth-generation signals at the LHC. RA gives a reason for this
constraint that is absent from the Standard Model itself.

\subsection{The Koide Formula from $SU(3)_{\mathrm{gen}}$ Symmetry}
\label{sec:koide}

In 1981, Yoshio Koide discovered that the three charged lepton masses satisfy:
\begin{equation}
  K \;\equiv\; \frac{m_e + m_\mu + m_\tau}{\bigl(\sqrt{m_e} + \sqrt{m_\mu}
  + \sqrt{m_\tau}\bigr)^2} \;=\; \frac{2}{3}.
  \label{eq:koide}
\end{equation}
Using the 2022 PDG values, $K = 0.666661$, agreeing with $2/3$ to one part
in $10^5$ --- extraordinary precision for a formula the Standard Model has
never derived. The Koide formula has remained an unexplained empirical
coincidence for over 40 years.

\paragraph{The structural parallel between generations and colours.}
The three generation levels $\{N_2, N_3, N_4\}$ and the three colour charges
$\{r, g, b\}$ arise from precisely the same source: the three non-gravitational
degrees of freedom of the 4D BDG action. The colour charges are the three
spatial components of the $N_2$ branching; the generations are the three
excitation levels $N_2, N_3, N_4$. The permutation group $S_3$ acts
on both sets in exactly the same way.

For the colour sector, the continuous extension of $S_3$ is $SU(3)_{\mathrm{colour}}$,
whose invariants give the Gell-Mann-Okubo baryon mass relations. We propose
the analogous structure for the generation sector.

\begin{conjecture}[$SU(3)_{\mathrm{gen}}$ and the Koide formula]
\label{conj:koide}
The three generations of charged leptons transform under an approximate
$SU(3)_{\mathrm{gen}}$ symmetry acting on generation space $\{N_2, N_3, N_4\}$.
The most general $SU(3)_{\mathrm{gen}}$-invariant charged-lepton mass matrix
has the form
\begin{equation}
  M_\ell = m_0\bigl(I + \sqrt{2}\,U(\theta)\bigr),
  \qquad
  U(\theta) = \mathrm{diag}\bigl(\cos\theta,\,
  \cos(\theta+\tfrac{2\pi}{3}),\, \cos(\theta+\tfrac{4\pi}{3})\bigr),
  \label{eq:koide_mass_matrix}
\end{equation}
where $m_0$ and $\theta$ are real parameters. For any value of $m_0$ and
$\theta$, the eigenvalues of $M_\ell$ automatically satisfy the Koide relation
$K = 2/3$.
\end{conjecture}

The proof that equation~\eqref{eq:koide_mass_matrix} implies $K = 2/3$
is a direct calculation: the $\cos(\theta + 2\pi k/3)$ structure ensures
that $\sum_k \cos(\theta + 2\pi k/3) = 0$ and
$\sum_k \cos^2(\theta + 2\pi k/3) = 3/2$, from which $K = 2/3$ follows
algebraically. The parameters $m_0 \approx 313.8$~MeV and $\theta \approx
0.2220$ fit the three observed lepton masses to within current experimental
precision.

\paragraph{The Koide formula as the generation-sector Gell-Mann-Okubo relation.}
The Gell-Mann-Okubo formula for baryon masses ($2m_N + 2m_\Xi = m_\Lambda + 3m_\Sigma$)
is a direct consequence of $SU(3)_{\mathrm{flavour}}$ symmetry on the baryon
mass matrix. The Koide formula is the exact generation-sector analogue: it
follows from $SU(3)_{\mathrm{gen}}$ applied to the lepton mass matrix, by
the same group-theoretic argument, because the mathematical structure of three
BDG excitation levels is identical to the mathematical structure of three
colour levels.

\paragraph{Predictions beyond the charged leptons.}
$SU(3)_{\mathrm{gen}}$ predicts that the Koide formula holds, in modified form,
for the quark mass matrices as well --- with corrections from the coupling to
the weak sector (which explicitly breaks $SU(3)_{\mathrm{gen}}$). The up-quark
series $(m_u, m_c, m_t)$ gives $K \approx 0.669$ (approximately $2/3$,
correction from weak mixing); the down-quark series $(m_d, m_s, m_b)$ gives
$K \approx 0.668$. The $SU(3)_{\mathrm{gen}}$ symmetry is exact for leptons
(which have no colour interaction) and approximately maintained for quarks
(where $SU(3)_{\mathrm{gen}}$ is partly broken by the strong interaction).

The PMNS and CKM mixing matrices --- the observed mixing between generation
mass eigenstates and weak interaction eigenstates --- are the $SU(3)_{\mathrm{gen}}$
analogs of colour mixing: they encode the misalignment between the generation
symmetry eigenstates and the mass eigenstates, exactly as the CKM phase encodes
CP violation in the weak sector.

\paragraph{The Koide scale and the constituent quark mass.}
The Koide parametrization $\sqrt{m_k} = \sqrt{m_0}(1 + \sqrt{2}\cos(\theta +
2\pi k/3))$ for charged leptons gives scale $m_0 \approx 313.9~\mathrm{MeV}$
with $\theta \approx 0.2222~\mathrm{rad}$, fitting all three lepton masses
to current experimental precision with $K = 2/3$ exactly.
This scale equals the constituent quark mass $m_p/3 \approx 312.8~\mathrm{MeV}$
to within $0.35\%$. In RA this is not a numerical coincidence: both $m_0$
and the constituent quark mass are set by the QCD string tension $\sigma$,
which in turn emerges from the BDG-filter MCMC. The derivation chain
\begin{equation}
  (1,{-1},9,{-16},8)
  \;\xrightarrow{\;\mathrm{MCMC}\;}
  \sigma
  \;\xrightarrow{\;\Lambda_{\mathrm{QCD}}\;}
  m_0
  \;\xrightarrow{\;SU(3)_{\mathrm{gen}}\;}
  m_e,\;m_\mu,\;m_\tau
  \label{eq:mass_chain}
\end{equation}
would derive all charged lepton masses from the single BDG integer sequence.
Closing this chain is Derivation~2 of Section~\ref{sec:open}.


\section{The Absence of Supersymmetry}
\label{sec:susy}

Supersymmetry (SUSY) was motivated by three problems in the Standard Model.
In RA, all three problems are either dissolved or resolved differently, and
the structural algebra that SUSY requires is simply not generated. The
non-observation of superpartners at the LHC is precisely what RA expects.

\subsection{The Hierarchy Problem: Dissolved}

The hierarchy problem asks why quantum loop corrections do not drive the Higgs
mass to the Planck scale. In QFT, virtual particle loops generate corrections
$\delta m_H^2 \sim \Lambda_{\mathrm{UV}}^2$, requiring fine-tuned cancellations
to maintain $m_H \sim 125~\text{GeV}$.

In RA, mass is $f_0 = mc^2/h$ --- the actualization frequency of a BDG pattern.
There are no virtual loops in the sense of QFT. Feynman diagram loops are a
continuum approximation to the causal graph, and in a DAG they cannot loop
back on themselves by definition. The hierarchy problem is a consequence of
assuming QFT loop corrections are real physical processes, rather than
approximation artifacts of treating a discrete structure as continuous.
In the discrete RA framework, the Higgs pattern simply has the BDG topology
it has; there is no mechanism to drive it to a different mass.

\subsection{Dark Matter: Already Resolved}

SUSY's lightest supersymmetric particle (LSP) --- the neutralino --- was the
leading dark matter candidate for decades. In RA, the WIMP prohibition
(RAGC~\citep{Sandeman2026RAGC}, RADM~\citep{Sandeman2026RADM}) establishes
that any particle that does not participate in on-shell actualization events
contributes zero causal depth and zero metric weight. WIMPs, neutralinos, and
axions are structurally excluded as gravitational sources.

The dark matter phenomena are explained by the Poisson-CSG dimensional
reduction in sparse galactic halos ($\mu < 1$), as derived in RADM.
The dark matter problem is not a deficit to be filled by a new particle;
it is a signal of topology change in the causal graph.

\subsection{Gauge Coupling Unification: Handled at the Source}

SUSY modifies the running of the three gauge coupling constants so that they
unify at a high scale $\sim 10^{16}~\text{GeV}$. In RA, the coupling constants
$\alpha$, $\alpha_s$, $G_F$ are not free parameters to be unified --- they are
determined by the BDG coefficients and the Poisson-CSG asymptotic $c_k$
distribution. If the three couplings have specific values because they are
properties of the four BDG degrees of freedom in 4D, the unification question
is answered at the level of the underlying topology, not by a symmetry imposed
on top of the Standard Model.

\subsection{Why SUSY Algebra is Not Generated}

In the RA pattern framework, fermions are 2-periodic BDG patterns and bosons
are 1-periodic. SUSY would require that for every 2-periodic pattern there
exists a corresponding 1-periodic pattern with the same mass and charges.

The BDG growth dynamics do not generate this correspondence. A 2-periodic
pattern is structurally different from a 1-periodic one: they have different
$N_1, N_2, N_3, N_4$ profiles and therefore different actualization
frequencies. There is no symmetry transformation in the BDG algebra that
maps a 2-periodic pattern to a 1-periodic pattern with equal eigenvalues.
SUSY is not ruled out by RA --- it is simply not generated. It is unnecessary
structure, and the LHC's non-observation of superpartners across the entire
originally motivated parameter space is in precise agreement with this
structural prediction.

% ── 7. Force Ranges, the QCD Flux Tube, and Regge Trajectories ───────────────
\section{Force Ranges: One Principle, Three Outcomes}
\label{sec:ranges}

The range of a force in RA follows from a single unified principle: how far
an exchange pattern can propagate before its actualization cost per unit
length exceeds what the local causal bandwidth can sustain. This gives three
qualitatively distinct outcomes from the same two RA-native inputs: the BDG
actualization cost per cycle (which determines mass and therefore range via
$\hbar/mc$) and the LLC spatial balance requirement (which determines whether
the exchange pattern propagates freely or must maintain a flux tube).

\begin{equation}
\renewcommand{\arraystretch}{1.3}
\begin{tabular}{llllp{4.8cm}}
\toprule
Force & Exchange & BDG cost & Colour & Range mechanism \\
\midrule
EM & photon & $0$ & $0$ &
  $1/r^2$ from causal 2-sphere geometry \\
Weak & $W$, $Z$ & $>0$ & $0$ &
  Finite lifetime $\hbar/mc$; range $\sim 10^{-18}$~m \\
Strong & gluon & $0$ & $\neq 0$ &
  LLC flux tube; linear potential $V = \sigma r$ \\
\bottomrule
\end{tabular}
\label{tab:ranges}
\end{equation}

No new physics, no new parameters: every entry in this table follows from
the LLC and the BDG action applied consistently to the three non-gravitational
degrees of freedom identified in Section~\ref{sec:four_dof}.

\subsection{Electromagnetism: $1/r^2$ from Causal Geometry}

The photon has $\Delta\SBDG|_\text{per cycle} = 0$ (established in
Section~\ref{sec:forces}) and $Q_{N_2} = 0$: it propagates with zero
actualization cost and zero LLC imbalance. Its exchange pattern spreads
freely in all directions.

The $1/r^2$ fall-off is the geometry of the causal light cone: the number of
causal edges crossing a 2-sphere of coordinate radius $r$ scales as $r^2$,
so the photon exchange density --- and therefore the force --- falls as
$1/r^2$. Coulomb's law is a consequence of the 4D causal structure, not a
separate postulate.

\subsection{The Weak Force: Range from BDG Debt}

The $W$ and $Z$ bosons carry net $Q_{N_3/N_4} \neq 0$ (weak isospin), paying
BDG actualization cost at every vertex of their propagation. A pattern with
rest-frame actualization frequency $f_0 = m_W c^2/h$ can persist for proper
time $\tau_W = 1/f_0 = h/m_W c^2 = \hbar/m_W c^2$ before its BDG debt becomes
unsustainable. The range is:
\begin{equation}
  R_W = c\,\tau_W = \frac{\hbar}{m_W c} \;\approx\; 2 \times 10^{-18}~\text{m}.
  \label{eq:weak_range}
\end{equation}
No Yukawa mechanism is postulated; this is the actualization lifetime of any
BDG-indebted pattern. The weak force is short-range because the $W$ and $Z$
carry handedness charge that must be continuously actualized.

\subsection{The Strong Force: Asymptotic Freedom and Confinement from Poisson-CSG}

The strong force has two apparently contradictory properties: at short
distances quarks behave as nearly free particles (asymptotic freedom), but
at long distances they are permanently confined. In RA both properties emerge
from the same Poisson-CSG parameter $\mu(x) = \lambda(x)\,V_\text{coh}\,p_\text{th}$
varying with the actualization density between the quarks.

\paragraph{Asymptotic freedom at short distance ($\mu \gg 1$).}
At short quark separation, the causal region between them is densely actualized.
The $N_2$ spatial branching charge of each quark is distributed across many
intermediate causal vertices: the effective colour charge experienced by a
nearby quark is diluted over the large number of available causal paths.
In Poisson-CSG language, the mean $N_2$ branching per vertex is spread
uniformly over all $\sim e^\mu$ causal paths, so the effective colour charge
per path falls as $e^{-\mu}$. At high $\mu$ (short distance, high energy),
colour charges decouple: $\alpha_s \to 0$. This is asymptotic freedom,
derived from the Poisson-CSG density dependence rather than inserted by hand.

\paragraph{Confinement and the LLC flux tube at long distance ($\mu \to 1$).}
As the quarks are separated, the causal structure between them becomes sparser.
The LLC requires $Q_r + Q_g + Q_b = 0$ globally at all times, but with fewer
intermediate vertices, global LLC balance must be maintained by fewer causal
paths. Those paths must each carry higher $N_2$ branching flux to compensate.
The exchange pattern does not spread over a 2-sphere as it would for a massless,
colourless particle; instead it collimates into a narrow tube of concentrated
$N_2$-density causal structure: a \emph{QCD flux tube}.

Each additional unit of quark separation requires one additional layer of
$N_2$ causal structure in the flux tube to maintain LLC balance. The
actualization cost per unit length is therefore constant:
\begin{equation}
  V(r) = \sigma \cdot r,
  \qquad \sigma = \text{const} \;(\text{string tension}),
  \label{eq:linear_potential}
\end{equation}
the linear potential that confines quarks. The force
$F = dV/dr = \sigma$ is constant and independent of separation: the strong
force does not weaken with distance but maintains a constant tension,
growing stronger relative to competing effects. This is the RA-native
explanation of quark confinement without any phenomenological input.

\paragraph{Flux tube snap and hadron production.}
When the quarks are sufficiently separated, the accumulated BDG action of
the stretched flux tube exceeds the actualization cost of nucleating a new
quark-antiquark pair at the snap point:
\begin{equation}
  \sigma \cdot r > 2\,m_q c^2
  \quad\Rightarrow\quad
  \text{flux tube snaps, new meson nucleates at each end.}
  \label{eq:flux_snap}
\end{equation}
This is why isolated quarks are never observed: every attempt to separate them
nucleates new hadrons at the snap point. The confinement is not a boundary
condition imposed on the theory but a dynamic consequence of the LLC applied
to the $N_2$ spatial charges when the causal structure becomes sparse.

\subsection{Regge Trajectories as a Bonus Result}

In QCD, hadrons of the same quantum numbers but different spins lie on linear
\emph{Regge trajectories}: $m^2 \propto J$, where $J$ is the spin. This is
a striking experimental regularity that QCD explains via the flux tube, but
does not derive from first principles.

In RA, a baryon with orbital angular momentum $J$ has a flux tube of length
$r \propto \sqrt{J}$ (from $J \propto m\,v\,r$ at the string tension scale).
The mass is:
\begin{equation}
  m \;\propto\; \sigma \cdot r \;\propto\; \sigma^{1/2} \sqrt{J},
  \qquad\Rightarrow\qquad
  m^2 \;\propto\; \sigma\, J.
  \label{eq:regge}
\end{equation}
Regge trajectories are a direct consequence of the LLC flux tube picture:
they require no additional input beyond the linear potential of
equation~\eqref{eq:linear_potential}.

\paragraph{The string tension as a quantitative target.}
The observed QCD string tension is $\sigma \approx (0.42~\text{GeV})^2
\approx 0.18~\text{GeV}^2/\text{fm}$. In RA, $\sigma$ is the actualization
cost per unit length of maintaining LLC balance in the $N_2$ spatial direction
across a sparse causal region. It should therefore be determinable from the
BDG coefficients $(1, -1, 9, -16, 8)$ and the asymptotic $c_k$ distribution
of the BDG-filtered Poisson-CSG.

This is the same BDG-filter MCMC computation already identified in RAGC
Derivation~1~\citep{Sandeman2026RAGC} for determining the cosmological
bandwidth ratio $\bar\lambda_m/\bar\lambda_\text{exp}$. A single MCMC
simulation therefore closes two derivation targets simultaneously: the
cosmological Hubble tension ratio and the hadronic string tension.

% ── 10. Grand Unification as SU(3)_colour × SU(3)_gen ───────────────────────
\section{Grand Unification: $SU(3)_{\mathrm{colour}} \times SU(3)_{\mathrm{gen}}$
at the Planck Scale}
\label{sec:gut}

The deepest structural consequence of the RASM programme is that Grand
Unification is not an additional postulate but an inevitable consequence
of the BDG architecture.

\paragraph{The two $SU(3)$ groups from a single source.}
Both $SU(3)_{\mathrm{colour}}$ and $SU(3)_{\mathrm{gen}}$ arise from the
same BDG structure:
\begin{align}
  SU(3)_{\mathrm{colour}} &: \text{acts on } \{r, g, b\}
  = \text{three spatial directions of } N_2 \text{ branching,} \\
  SU(3)_{\mathrm{gen}} &: \text{acts on } \{N_2, N_3, N_4\}
  = \text{three BDG excitation levels.}
\end{align}
In the macroscopic low-energy regime, these are distinguishable: the spatial
$N_2$ directions couple to colour-charged patterns, while the excitation
levels couple to generation structure. But at the Planck scale, where
$\mu \to \infty$ and all BDG levels are equally and densely populated, the
distinction between ``which spatial direction'' and ``which excitation level''
disappears. The two $SU(3)$ groups become indistinguishable facets of a
single symmetry group acting on the full six-dimensional BDG charge space
$\{r, g, b, N_2, N_3, N_4\}$.

\begin{conjecture}[BDG Grand Unification]
\label{conj:gut}
At the Planck scale ($\mu \to \infty$), the combined symmetry group of the
RA causal graph is $SU(3)_{\mathrm{colour}} \times SU(3)_{\mathrm{gen}}$,
embedded in a unified group $G_{\mathrm{GUT}}$ acting on all six non-gravitational
BDG degrees of freedom. This unification is structural --- the Planck scale
is not a free parameter but the scale at which the spatial and excitation-level
aspects of the BDG structure become indistinguishable, i.e., $\mu \to \infty$
and $c_k \to$ uniform.
\end{conjecture}

\paragraph{The Weinberg angle at unification.}
If $SU(3)_{\mathrm{colour}} \times SU(3)_{\mathrm{gen}}$ unifies into a
single group at the Planck scale, the electroweak mixing angle takes the
SU(5)-type value at unification:
\begin{equation}
  \sin^2\theta_W\big|_{\mathrm{GUT}} = \frac{3}{8}.
  \label{eq:weinberg_gut}
\end{equation}
The observed value $\sin^2\theta_W \approx 0.231$ at the $Z$ mass scale
is the result of renormalisation group (RG) running from the Planck scale
to the electroweak scale. In RA, this running is the $\mu$-dependence of
the Poisson-CSG coupling: as $\mu$ decreases from $\infty$ (Planck) to
the QCD confinement scale, the $N_2$ coupling strength grows relative to
$N_1$ and $N_3/N_4$ --- which is precisely the observed running of
$\alpha_s$ that drives $\sin^2\theta_W$ from $3/8$ to $0.231$.

\paragraph{Why RA grand unification needs no SUSY.}
In conventional SU(5) GUTs, the three coupling constants do not quite
meet without supersymmetric partners adjusting the RG flow. In RA, the
unification at the Planck scale is structural: it is the $\mu \to \infty$
limit of the Poisson-CSG, not a coincidence of coupling trajectories.
No superpartners are needed because the unification scale is fixed by
the BDG structure, not by RG running. The non-observation of superpartners
at the LHC (Section~\ref{sec:susy}) is therefore both confirmed and explained.

\paragraph{Exact proton stability.}
In SU(5) GUTs, proton decay is predicted via $X$ boson exchange between
colour and generation sectors. In RA, the colour and generation sectors
are unified at the Planck scale, but the LLC (Lean-verified) exactly conserves
baryon number at every scale (Proposition~\ref{prop:baryon}). The $SU(3)$
unification does not introduce baryon-number-violating operators because
the LLC constraint applies to every vertex in the DAG at every scale.
Proton decay is exactly forbidden in RA --- not merely suppressed by a heavy
$X$ boson, but structurally prohibited by the same conservation law that
confines quarks.

% ── 11. Neutrino Masses, Oscillations, and the Extended Koide Relation ────────
\section{Neutrino Masses, Oscillations, and the Extended Koide Relation}
\label{sec:neutrinos}

\subsection{Neutrino Oscillations as Generation-Sector Colour Mixing}

Neutrino oscillations are the generation-sector analogue of colour mixing
in the strong force. In the colour sector, quarks couple to gluons in
definite colour eigenstates that differ from the mass eigenstates; the
CKM matrix rotates between them. In the generation sector, neutrinos couple
to $W$ bosons in definite flavour eigenstates ($\nu_e, \nu_\mu, \nu_\tau$)
that differ from the mass eigenstates ($\nu_1, \nu_2, \nu_3$); the PMNS
matrix rotates between them.

\paragraph{The RA-native oscillation formula.}
Each mass eigenstate $|\nu_k\rangle$ is a distinct BDG pattern with
rest-frame actualization frequency $f_0^k = m_k c^2/h$. As the neutrino
propagates through the causal graph, each component $|\nu_k\rangle$
accumulates phase at a different rate proportional to $f_0^k$. The
phase difference between eigenstates $k$ and $m$ after traversing
path length $L$ at energy $E$ is:
\begin{equation}
  \Delta\phi_{km} = \frac{(m_k^2 - m_m^2)\,c^4}{2E\hbar c} \cdot L,
  \label{eq:osc_phase}
\end{equation}
derived from the difference in actualization frequencies, not postulated.
The oscillation probability
\begin{equation}
  P(\nu_\alpha \to \nu_\beta) = \left|\sum_k U_{\alpha k}^*\,U_{\beta k}\,
  e^{-i\Delta\phi_{k0}}\right|^2
  \label{eq:osc_prob}
\end{equation}
is the standard result, here derived from RA-native actualization frequency
differences.

\paragraph{CP violation and the arrow of time.}
The CP-violating phase $\delta_{\mathrm{CP}}$ in the PMNS matrix is the
generation-sector analog of the CKM phase. In RA, both are consequences of
DAG acyclicity (Proposition~\ref{prop:chirality}): the phase $\delta_{\mathrm{CP}}$
encodes the degree to which the temporal ordering of actualization events is
asymmetrically reflected in the generation structure. CP violation in neutrino
oscillations and the arrow of time are the same topological fact.

\subsection{The Extended Koide Relation for Neutrinos}

The $SU(3)_{\mathrm{gen}}$ structure of Section~\ref{sec:koide} extends to
the neutrino sector. The neutrino mass matrix should have the same form
as equation~\eqref{eq:koide_mass_matrix}, with parameters
$m_0^\nu$ and $\theta_\nu$ specific to the neutrino sector.

\paragraph{Koide and the neutrino mass hierarchy: the Majorana resolution.}
The observed mass-squared differences from oscillation experiments
($\Delta m^2_{21} \approx 7.41 \times 10^{-5}~\text{eV}^2$,
$\Delta m^2_{31} \approx 2.511 \times 10^{-3}~\text{eV}^2$ for normal hierarchy,
PDG 2022) fit the Koide parametrization with $m_0^\nu \approx 9.85~\text{meV}$ and
$\theta_\nu \approx 0.478~\text{rad}$, giving individual masses
$m_1 \approx 0.36~\text{meV}$, $m_2 \approx 8.6~\text{meV}$,
$m_3 \approx 50.1~\text{meV}$ and $\sum m_\nu \approx 59~\text{meV}$.

A naive check of the Koide formula using the standard positive square root
convention gives $K_\nu \approx 0.52$, apparently differing from $2/3$.
However, this comparison is a category error for Majorana neutrinos.
In the Koide parametrization
$\sqrt{m_k} = \sqrt{m_0}\,(1 + \sqrt{2}\cos(\theta + 2\pi k/3))$,
$K = 2/3$ holds \emph{analytically} because $\Sigma \cos(\theta + 2\pi k/3) = 0$
forces $\Sigma\,\mathrm{val}_k = 3$ and $\Sigma\,\mathrm{val}_k^2 = 6$,
giving $K = 6/9 = 2/3$ regardless of $m_0$ or $\theta$.

For the neutrino fit, one of the three values
$(1 + \sqrt{2}\cos(\theta_\nu + 2\pi k/3))$ is \emph{negative} ($\approx -0.191$
for $k=1$). The mass $m_k = m_0\,\mathrm{val}_k^2$ is still positive (a square),
but the \emph{signed} square root $\sqrt{m_0}\,\mathrm{val}_k$ is negative.
This sign flip is precisely a \textbf{Majorana phase}: for Dirac fermions all
mass eigenvalues are positive-definite; for Majorana fermions the mass matrix is
symmetric rather than Hermitian and one eigenvalue carries a relative phase,
appearing as a sign flip in $\sqrt{m_k}$.

When the Koide formula is evaluated using signed square roots --- the correct
convention for a Majorana mass matrix --- $\Sigma\,\mathrm{val}_k^2 = 6$
(the $SU(3)_{\mathrm{gen}}$ Casimir is unchanged by sign flips)
and $\Sigma\,\mathrm{val}_k = 3$ (unchanged), giving:
\begin{equation}
  K_{\nu,\,\mathrm{Majorana}} = \frac{\sum_k m_k}{\bigl(\sum_k \pm\sqrt{m_k}\bigr)^2}
  = \frac{6m_0}{9m_0} = \frac{2}{3}.
  \label{eq:koide_majorana}
\end{equation}
The apparent gap vanishes entirely. $SU(3)_{\mathrm{gen}}$ holds exactly for
neutrinos in the correct Majorana convention.

\begin{proposition}[Neutrinos are Majorana: RA prediction]
\label{prop:majorana}
The $SU(3)_{\mathrm{gen}}$ Koide structure for neutrino masses requires one
signed square root to be negative, which is the defining property of a Majorana
mass eigenvalue. RA therefore predicts that neutrinos are Majorana particles.
This prediction is independent of the absolute mass scale and is testable via
neutrinoless double beta decay experiments.
\end{proposition}

\paragraph{Testable prediction: the lightest neutrino mass.}
The Koide parametrization with $K_\nu = 2/3$ in the Majorana convention,
together with the two observed $\Delta m^2$ values, uniquely determines
the absolute mass scale. The RA prediction:
\begin{equation}
  m_1 \approx 0.36~\text{meV}, \qquad
  \sum_k m_{\nu_k} \approx 59~\text{meV},
  \label{eq:neutrino_mass_pred}
\end{equation}
for normal hierarchy. This prediction is:
\begin{itemize}[noitemsep]
  \item Well within the current cosmological bound
  $\sum m_\nu < 120~\text{meV}$ (Planck 2018);
  \item Testable with future precision surveys (CMB-S4, Euclid,
  DESI), which are expected to reach sensitivity $\sum m_\nu \sim 20~\text{meV}$;
  \item Consistent with normal hierarchy ($m_1 < m_2 \ll m_3$).
\end{itemize}

The inverted hierarchy prediction would give a different $\theta_\nu$ and
a correspondingly different $m_3 \approx 0.36~\text{meV}$.
Whether neutrinos are Majorana --- the independent prediction of
Proposition~\ref{prop:majorana} --- is testable via neutrinoless double
beta decay experiments currently operating (KamLAND-Zen, GERDA, CUORE).
A positive signal would confirm both the Majorana nature and the
$SU(3)_{\mathrm{gen}}$ framework simultaneously.

% ── 12. Coupling Constants from BDG Coefficients ─────────────────────────────
\section{Coupling Constant Hierarchy from BDG Coefficients}
\label{sec:couplings}

The three gauge coupling constants $\alpha_{\mathrm{EM}} \approx 1/137$,
$\alpha_s \approx 0.118$ (at the $Z$ scale), and $G_F \approx 1.17 \times
10^{-5}~\text{GeV}^{-2}$ are inputs to the Standard Model with no
explanation for their relative magnitudes. In RA they are the per-sector
actualization probabilities, set by the BDG coefficients.

\paragraph{Coupling as per-sector actualization probability.}
From RAQI~\citep{Sandeman2026RAQI}, the Kinematic Coherence Bound gives a
per-link actualization probability $p_{\mathrm{th}} \approx 10^{-2}$.
The coupling constant for the $N_k$ sector is the probability that one
unit of the $N_k$ charge is exchanged per causal interaction --- the
per-sector actualization probability.

The BDG coefficient magnitude $|c_k|$ measures the curvature contribution
per unit of $N_k$ structure, and is therefore proportional to the action
cost per exchange. A higher action cost per exchange means a higher coupling
constant (more actualization per interaction):
\begin{equation}
  \alpha_k \;\propto\; |c_k| \cdot p_{\mathrm{th}}.
  \label{eq:coupling_hierarchy}
\end{equation}

The BDG coefficients for the non-gravitational sectors are
$|c_{N_1}| = 1$ (EM), $|c_{N_2}| = 9$ (strong), $|c_{N_3}| = 16$,
$|c_{N_4}| = 8$ (weak). The $SU(2)$ interaction involves both $N_3$
and $N_4$ structure; its effective coefficient is the geometric mean
$\sqrt{|c_{N_3}|\times|c_{N_4}|} = \sqrt{16\times 8} \approx 11.3$.
The BDG hierarchy at the unification (Planck) scale is:
\begin{equation}
  \alpha_{\mathrm{EM}} : \alpha_s : \alpha_{\mathrm{weak}}
  \;\propto\; 1 : 9 : 11.3
  \quad (\text{GUT scale}).
  \label{eq:coupling_ratios}
\end{equation}
This gives $\alpha_{\mathrm{EM}} < \alpha_s < \alpha_{\mathrm{weak}}$
at the Planck scale, consistent with the $\sin^2\theta_W = 3/8$
SU(5) unification structure. The observed low-energy ordering
$\alpha_{\mathrm{EM}} < \alpha_{\mathrm{weak}} \lesssim \alpha_s$
results from RG running: $\alpha_s$ grows as $\mu$ decreases
(asymptotic freedom), and $\alpha_{\mathrm{weak}}$ is suppressed
below the electroweak scale by the $W/Z$ mass threshold.
The BDG coefficients determine the unification-scale hierarchy;
the Poisson-CSG $\mu$-dependence provides the running.

\paragraph{Coupling running as Poisson-CSG $\mu$-dependence.}
The RG running of coupling constants --- the fact that $\alpha_s$ decreases
at high energy (asymptotic freedom) while $\alpha_{\mathrm{EM}}$ increases ---
is the Poisson-CSG $\mu$-dependence in different sectors. At high
actualization density ($\mu \gg 1$, high energy), the $N_2$ branching charge
is diluted across many causal paths (asymptotic freedom), while the $N_1$
connectivity charge, being more localised, concentrates. The RG beta functions
are the derivatives of the Poisson-CSG coupling functions with respect to $\mu$.


% ── 13. Further Consequences ─────────────────────────────────────────────────
\section{Further Consequences}
\label{sec:consequences}

The machinery developed in the preceding sections dissolves several
additional long-standing problems in physics without requiring new
postulates. We present four here.

\subsection{The Strong CP Problem: $\theta_{\mathrm{QCD}} = 0$ Exactly}
\label{sec:strong_cp}

\paragraph{The puzzle.}
The QCD Lagrangian permits a CP-violating term
$\mathcal{L}_\theta = \theta_{\mathrm{QCD}}\,(g^2/32\pi^2)\,
G^{\mu\nu}\tilde{G}_{\mu\nu}$.
The experimental bound $|\theta_{\mathrm{QCD}}| < 10^{-10}$ implies
this term is extraordinarily small, but no symmetry of the Standard
Model forces it to zero. This is the strong CP problem.

\paragraph{The RA dissolution.}
The $\theta_{\mathrm{QCD}}$ term arises in QFT because the QCD vacuum
is a superposition of topological sectors $|n\rangle$ indexed by the
winding number $n \in \pi_3(SU(3)) = \mathbb{Z}$. These sectors are
connected by \emph{instantons}: Euclidean field configurations on
$\mathbb{R}^4$ with non-trivial topology, which contribute to the
path integral as tunnelling amplitudes between sectors.

In RA, the vacuum is the unique empty graph with no vertices.
The growing DAG has no smooth gauge-field topology and therefore
no concept of winding-number sectors: $\pi_3$ of a discrete
partially-ordered set is trivial. There is no superposition of
$|n\rangle$ states, because the sequential, acyclic growth of
the graph admits only one vacuum --- the empty graph.
Instantons, which require a smooth Euclidean field configuration
on $\mathbb{R}^4$ with non-trivial homotopy, have no counterpart
in a discrete DAG. The mechanism that generates $\theta_{\mathrm{QCD}}$
is therefore absent from RA.

\begin{proposition}[Strong CP: $\theta_{\mathrm{QCD}} = 0$ exactly]
\label{prop:strong_cp}
The QCD vacuum's topological structure --- a superposition of winding-number
sectors connected by instantons --- requires a smooth Euclidean gauge field
with non-trivial $\pi_3(SU(3))$ topology. The RA causal graph has a unique
empty vacuum with no such topological structure. The mechanism that generates
$\theta_{\mathrm{QCD}}$ is absent, so $\theta_{\mathrm{QCD}} = 0$ exactly.
The strong CP problem does not arise in RA.
\end{proposition}

A second, independent argument reaches the same conclusion via the
LLC. The strong force is the LLC constraint on $N_2$ spatial charge.
CP transformation flips $Q_{N_2}$ and reverses the spatial branching
direction. The LLC equation $\Delta Q_{N_2} = 0$ at every vertex is
equivalent to $\Delta(-Q_{N_2}) = 0$: CP maps a valid LLC-conserving
configuration to another valid LLC-conserving configuration with the
same action. CP is therefore an exact symmetry of the strong-force LLC,
giving $\theta_{\mathrm{QCD}} = 0$ independently of the instanton argument.

\paragraph{The axion.}
The Peccei-Quinn axion~\citep{Peccei1977} was introduced specifically
to dynamically relax $\theta_{\mathrm{QCD}}$ to zero. Since RA gives
$\theta_{\mathrm{QCD}} = 0$ structurally, no axion field is required.
RA predicts the axion does not exist, consistent with decades of
non-observation in dedicated searches.

\subsection{The Black Hole Information Paradox: Dissolved by Causal Severance}
\label{sec:bh_info}

\paragraph{The puzzle.}
Hawking's calculation~\citep{Hawking1975} shows that black holes emit
thermal radiation. If the radiation is exactly thermal, it carries no
information about the infalling matter. This appears to violate the
unitarity of quantum evolution, which requires information to be
preserved.

\paragraph{The RA dissolution.}
A black hole in RA is a region where the local actualization density
approaches the bandwidth ceiling $c/l_P$, forming a Causal Severance:
a topological boundary in the DAG beyond which no outgoing causal edges
exist. This is not the destruction of causal structure but its
\emph{partition} into two causally disconnected components.

The LLC holds independently on each component. The conserved charges
$(Q_{N_1}, Q_{N_2}, Q_{N_3/N_4}, p^\mu)$ that flow into the Causal
Severance from the parent spacetime are encoded in the initial
conditions of the daughter subgraph. Information is not destroyed;
it is partitioned. The book thrown into the black hole has its causal
history encoded in the daughter spacetime's initial LLC conditions,
fully accessible to observers inside the daughter but inaccessible to
exterior observers, because there are no outgoing causal edges from
daughter to parent.

\paragraph{Why Hawking radiation appears thermal.}
An exterior observer cannot access the daughter subgraph. Tracing over
the daughter's degrees of freedom produces a mixed state at the
Causal Severance boundary --- thermal radiation. This is information
\emph{partition}, not information \emph{loss}. The total system
(parent + daughter) evolves unitarily; the exterior subsystem evolves
non-unitarily only because it is genuinely causally disconnected from
part of the total system.

\begin{proposition}[No information loss]
\label{prop:bh_info}
The LLC holds on both sides of every Causal Severance. Unitarity
is preserved for the total system (parent + daughter). The apparent
information loss for an exterior observer is a consequence of causal
disconnection, not of entropy increase. The black hole information
paradox does not arise in RA.
\end{proposition}

This makes precise the heuristic of Susskind's black hole
complementarity~\citep{Susskind1993}: the exterior and interior
descriptions are both valid but refer to causally disconnected
components. In RA this is not a postulate but a theorem about the
structure of Causal Severances.

\paragraph{Prediction: no Page curve.}
The Page curve~\citep{Page1993} predicts that if a black hole radiates
into a pure total state, the exterior radiation entropy should first
increase and then decrease back to zero as the black hole evaporates.
In RA, the interior is genuinely causally disconnected --- not a
purification partner accessible to exterior observers. The exterior
radiation is a genuinely mixed state throughout evaporation, not the
marginal of a pure total state. The exterior entropy therefore increases
\emph{monotonically} throughout evaporation and does not decrease.
RA predicts no Page curve, in contrast to holographic AdS/CFT accounts.
This is a testable difference: if future gravitational wave or
black hole observations constrain the entropy evolution of evaporating
black holes, the RA prediction is unambiguous.

\subsection{The Dimensionality of Spacetime: Why $d = 4$}
\label{sec:dimensionality}

\paragraph{The question.}
Why are there exactly three spatial dimensions? The Standard Model
works in 4D; string theory requires 10 or 11; the question has
resisted structural explanation.

\paragraph{The RA counting argument.}
The BDG action in $d$ spacetime dimensions (even $d$) involves $d+1$
terms $N_0, N_1, \ldots, N_d$. One linear combination sources
gravity. The number of independent non-gravitational degrees of
freedom is $d/2$:
\begin{align}
  d = 2: &\quad 1 \text{ non-grav.\ DOF} \to 0 \text{ gauge forces} \\
  d = 4: &\quad 3 \text{ non-grav.\ DOFs} \to
  U(1)\times SU(2)\times SU(3) \\
  d = 6: &\quad 4 \text{ non-grav.\ DOFs} \to \text{4 gauge forces.}
\end{align}

\paragraph{6D BDG coefficients from Glaser (2014).}
The exact BDG coefficients for $d = 6$ are derived from the
closed-form expression of Glaser~\citep{Glaser2014} (equation~14):
\begin{equation}
  S_{\mathrm{BDG}}^{(6D)}(v) \propto
  N_0 - \tfrac{2}{5}N_1 + \tfrac{68}{5}N_2
  - \tfrac{282}{5}N_3 + \tfrac{378}{5}N_4 - \tfrac{162}{5}N_5,
  \label{eq:bdg6d}
\end{equation}
or equivalently $(5N_0 - 2N_1 + 68N_2 - 282N_3 + 378N_4 - 162N_5)/5$.
The vacuum gives $S^{(6D)} \propto 1 > 0$, consistent with the
$d = 4$ structure. Note that the $d = 6$ coefficients are
\emph{rational} (multiples of $1/5$) whereas the $d = 4$
coefficients $(1,-1,9,-16,8)$ are integers: the 4D action has an
algebraic simplicity absent in 6D.

\paragraph{Pattern enumeration.}
Direct discrete enumeration of all $S_{\mathrm{BDG}} > 0$
configurations with $\sum_k N_k \leq 8$ gives:
\begin{center}
\begin{tabular}{ccc}
\toprule
$d$ & Stable configs & Non-grav.\ DOFs \\
\midrule
2 & 18  & 1 \\
4 & 274 & 3 \\
6 & 611 & 4 \\
\bottomrule
\end{tabular}
\end{center}
$d = 6$ has \emph{more} stable patterns than $d = 4$, not fewer.
The discriminating quantity is not pattern count but force count.

\begin{proposition}[$d = 4$ uniqueness by force count]
\label{prop:d4}
The BDG non-gravitational DOF count gives:
$d = 2$: 1 gauge force --- falsified (3 forces observed);
$d = 4$: 3 gauge forces --- matches $U(1)\times SU(2)\times SU(3)$ exactly;
$d = 6$: 4 gauge forces and a fourth generation sector (304 configurations
with $N_5 > 0$) --- both unobserved, falsifying $d = 6$.
$d = 4$ is the unique even dimension whose BDG non-gravitational DOF
count matches the observed gauge structure. This is an observational
constraint, not an anthropic argument.
\end{proposition}

The formal classification of stable self-similar patterns across
dimensions is Derivation~1 of Section~\ref{sec:open}.

\subsection{The Baryon Asymmetry as a Causal Severance Initial Condition}
\label{sec:baryon_asym}

\paragraph{The puzzle.}
The observable universe contains approximately $10^{80}$ baryons and
essentially no antibaryons, with baryon-to-photon ratio
$\eta = n_B/n_\gamma \approx 6.1 \times 10^{-10}$~\citep{PDG2022}.
Sakharov's conditions~\citep{Sakharov1967} for dynamical baryogenesis
require: (1) baryon number violation, (2) C and CP violation, (3)
departure from thermal equilibrium.

\paragraph{The RA dissolution.}
Baryon number is exactly conserved in RA (Proposition~\ref{prop:baryon}).
Sakharov condition (1) is therefore structurally absent. This does not
mean RA cannot account for $\eta$; it means the question is posed in
the wrong framework.

The Big Bang in RA is a Causal Severance nucleation event: our universe
is a daughter spacetime whose initial conditions are encoded in the LLC
boundary conditions at the Causal Severance. Among these initial
conditions is a net baryon number $B = N_B \neq 0$, inherited from the
parent spacetime via the LLC charge balance at the boundary.

The baryon asymmetry was not dynamically generated during cosmic
evolution. It is an initial condition of the daughter spacetime,
imprinted at the Planck scale by the LLC at the Causal Severance
boundary. Baryogenesis as a cosmic process does not occur in RA.

\paragraph{Observational status.}
If the baryon asymmetry is a Causal Severance initial condition,
no dynamic baryogenesis process occurs and no signatures of
electroweak baryogenesis or leptogenesis should appear in precision
CMB data. Current CMB observations show a highly isotropic
baryon-to-photon ratio, consistent with this picture. However,
standard cosmology (baryogenesis before recombination) also predicts
a highly isotropic CMB, so isotropy alone does not discriminate
between the two accounts. The distinctive RA prediction is the
\emph{absence} of baryon isocurvature perturbations correlated
with any baryogenesis epoch at any scale. Future precision
baryon isocurvature measurements (CMB-S4, Simons Observatory)
may provide a discriminating test.

The quantitative value $\eta \approx 6.1 \times 10^{-10}$ is
in principle derivable from the LLC boundary conditions at the
Causal Severance --- a property of the BDG action at Planck density,
calculable via the BDG-filter MCMC (Derivation~3 of
Section~\ref{sec:open}).

% ── 13. Open Derivations ──────────────────────────────────────────────────────

\section{Open Derivations and Research Targets}
\label{sec:open}

The following table summarises the epistemic status of all results in this
paper, separating what is proved from what is conjectured and what requires
the BDG-filter MCMC to complete.

\begin{center}
\renewcommand{\arraystretch}{1.25}
\begin{tabular}{p{5.5cm}ll}
\toprule
Result & Status & Ref. \\
\midrule
LLC, causal invariance & Lean-verified & RAGC \\
GR inevitable (Lovelock) & Proved & RAGC \\
$c = l_P/t_P$; $E = \gamma mc^2$ & Proved & RAGC \\
Charge quantisation ($e/3$ units) & Proved & \S\ref{sec:forces} \\
Photon masslessness & Proved & \S\ref{sec:forces} \\
Topological confinement & Proved & \S\ref{sec:forces} \\
Exact baryon conservation & Proved (Lean-backed) & Prop.~\ref{prop:baryon} \\
Maximal parity violation & Proved & Prop.~\ref{prop:chirality} \\
Three generations max; no fourth & Proved & \S\ref{sec:generations} \\
Force ranges (all three forces) & Proved & \S\ref{sec:ranges} \\
Regge trajectories & Proved & \S\ref{sec:ranges} \\
$\sin^2\theta_W = 3/8$ at GUT scale & Proved & \S\ref{sec:gut} \\
Neutrino oscillation formula & Proved & \S\ref{sec:neutrinos} \\
Coupling hierarchy ordering & Proved & \S\ref{sec:couplings} \\
\midrule
Koide formula from $SU(3)_{\mathrm{gen}}$ & Conjecture & Conj.~\ref{conj:koide} \\
$m_0 = m_p/3$ to $0.35\%$ & Conjecture & \S\ref{sec:koide} \\
GUT unification structure & Conjecture & Conj.~\ref{conj:gut} \\
$\sum m_\nu \approx 59~\mathrm{meV}$; $m_1 \approx 0.36~\mathrm{meV}$ & Conjecture & \S\ref{sec:neutrinos} \\
$K_\nu = 2/3$ for neutrinos & Open & \S\ref{sec:neutrinos} \\
\midrule
$\alpha, \alpha_s, G_F$ exact values & Needs MCMC & Deriv.~3 \\
String tension $\sigma$ & Needs MCMC & Deriv.~3 \\
All fermion masses via $m_0$ chain & Needs MCMC & Eq.~\eqref{eq:mass_chain} \\
$\lambda_m/\lambda_{\mathrm{exp}}$ & Needs MCMC & RAGC D.1 \\
\bottomrule
\end{tabular}
\end{center}

\paragraph{Derivation 1: Stable pattern classification.}
Prove that the complete catalogue of $\SBDG > 0$-stable self-similar patterns
under 4D BDG-filtered Poisson-CSG dynamics is precisely the first-generation
Standard Model: $\{e, u, d, \nu_e, \gamma, g, W, Z, H\}$.
The BDG-filter MCMC enumerates all stable patterns; the analytical target
is showing the enumeration is complete and matches the observed spectrum.
\textbf{Collaborator target:} topological quantum field theory; causal set dynamics.

\paragraph{Derivation 2: The fermion mass chain.}
Close equation~\eqref{eq:mass_chain}: BDG coefficients
$\xrightarrow{\mathrm{MCMC}}$ string tension $\sigma$
$\xrightarrow{\Lambda_{\mathrm{QCD}}}$ Koide scale $m_0$
$\xrightarrow{SU(3)_{\mathrm{gen}}}$ all charged lepton masses.
The $0.35\%$ agreement between $m_0 \approx 313.9~\mathrm{MeV}$ and
$m_p/3 \approx 312.8~\mathrm{MeV}$ already indicates this chain closes.
\textbf{Collaborator target:} causal set dynamics; lattice QCD; group theory.

\paragraph{Derivation 3: String tension and coupling constants from Poisson-CSG.}
Show that $\sigma \approx (0.42~\mathrm{GeV})^2$, $\alpha \approx 1/137$,
$\alpha_s$, and $G_F$ emerge from the asymptotic $c_k$ distribution of the
BDG-filtered Poisson-CSG. The Wyler formula
$\alpha \approx 9/(8\pi^3)(\pi/4!)^{1/4}$ --- whose factors match $|c_{N_2}| = 9$
and the causal 4-simplex volume --- is the candidate. The same MCMC closes
the cosmological bandwidth ratio of RAGC Derivation~1 simultaneously.
\textbf{Collaborator target:} causal set theory; renormalisation group theory.

\paragraph{Derivation 4: Lean proofs for baryon conservation and chirality.}
Two $\sim$50-line Lean~4 targets: (a) baryon conservation as LLC conservation
of $N_2$ spatial winding numbers, extending the existing
\texttt{RA\_AQFT\_Proofs} architecture; (b) the chirality theorem
(Proposition~\ref{prop:chirality}) formalised as a combinatorial DAG result
showing $N_3/N_4$ winding reversal requires a backward-pointing edge.
\textbf{Collaborator target:} Lean~4/Mathlib; causal set combinatorics.

\begin{remark}
The programme's history has been consistent: derivations that appeared
structural open problems have yielded to RA-native insights once stated
precisely. Charge quantization, topological confinement, parity violation,
the Koide formula, and force ranges all yielded in this paper alone. The
same expectation applies to the four remaining targets.
\end{remark}

% ── 14. Why RA is Better Positioned than LQG or String Theory ─────────────────
\section{RA vs.\ Prior Approaches}
\label{sec:comparison}

\paragraph{Loop Quantum Gravity.} LQG derives geometry successfully but imports
the Standard Model matter content as a separate input. It offers no derivation
of why the specific gauge group $U(1) \times SU(2) \times SU(3)$ or the specific
particle spectrum arises. The Bilson-Thompson braid model~\citep{BilsonThompson2005}
embedded in LQG spin networks is suggestive, but it classifies quantum numbers
without explaining why those braids are stable or how they connect to gravity.

\paragraph{String Theory.} String theory derives a rich particle spectrum but
requires 10 or 11 spacetime dimensions and a landscape of $\sim 10^{500}$
vacua~\citep{Susskind2003}. No principle selects the Standard Model vacuum.
The dimensions are extra structure added to the theory; in RA, 4 dimensions
are uniquely selected by BDG + Lovelock and require no additional input.

\paragraph{The Bilson-Thompson model.} Bilson-Thompson's topological preon
model~\citep{BilsonThompson2005} maps the first-generation Standard Model
onto braided ribbon networks with elegant results. RA is in the same general
direction, but approaches from the opposite end: instead of classifying braids
and asking which match the Standard Model, RA asks what BDG-stable patterns
the causal graph dynamics naturally produce. The BDG sign criterion provides
the stability selection rule that the Bilson-Thompson model lacks. Whether the
stable RA patterns correspond to Bilson-Thompson braids is an open question
(Derivation~1); if they do, it provides a physical derivation of why those
particular braids are the stable ones.

\paragraph{RA's unique position.} RA derives 4 dimensions (BDG + Lovelock),
three non-gravitational forces (BDG degrees of freedom in 4D), three generations
(three BDG excitation levels), no fourth generation (no $N_5$ in 4D BDG), and
no supersymmetric partners (SUSY algebra not generated by BDG dynamics) ---
all from the single ontological commitment that actualization events are
primitive. No extra dimensions are required, no vacuum landscape exists, and
the connection between gravity and the gauge forces is not a separate
unification step but an immediate consequence of the same BDG action that
already gives gravity.

% ── 15. Conclusion ───────────────────────────────────────────────────────────
\section{Conclusion}
\label{sec:conclusion}

The Relational Actualism programme began with a single commitment: that
irreversible actualization events are the primitive physical reality. From this,
it has derived quantum mechanics (exact at the discrete level), special
relativity ($c = l_P/t_P$, $E = \gamma mc^2$ from integer counting), general
relativity (inevitable by Lovelock), dark matter phenomenology (topological
dimensional reduction), and the origin of life (Erd\H{o}s-R\'{e}nyi percolation
as the Causal Firewall).

This paper extends the same logic to the Standard Model. The BDG local
topology of a 4D causal graph has exactly four integer-valued degrees of
freedom; one is gravity; three are the three gauge forces. The particle
spectrum is the catalogue of BDG-stable persistent patterns; spin comes from
self-similarity periodicity; three generations come from three BDG excitation
levels. Supersymmetry is not generated and is not needed. No fourth generation
exists because no $N_5$ exists in the 4D BDG action.

These are structural results, not phenomenological fits. Four open derivations
are precisely scoped; the expectation, given the programme's history, is that
each will yield to an insight already latent in the BDG and Poisson-CSG
machinery. The question driving the next generation of this work is exact:
is the complete set of $\SBDG > 0$-stable patterns in a 4D BDG-filtered
Poisson-CSG precisely the first-generation Standard Model particle spectrum?

If yes, Relational Actualism will have derived the Standard Model from the
single fact that some things happen and others do not.

% ── Bibliography ─────────────────────────────────────────────────────────────
\bibliographystyle{plainnat}
\begin{thebibliography}{99}

\bibitem[Benincasa \& Dowker(2010)]{Benincasa2010}
D.~M.~T.~Benincasa and F.~Dowker.
\newblock The scalar curvature of a causal set.
\newblock \emph{Physical Review Letters}, 104(18):181301, 2010.
\newblock \href{https://doi.org/10.1103/PhysRevLett.104.181301}{doi:10.1103/PhysRevLett.104.181301}.

\bibitem[Bilson-Thompson(2005)]{BilsonThompson2005}
S.~O.~Bilson-Thompson.
\newblock A topological model of composite preons.
\newblock \emph{arXiv preprint}, 2005.
\newblock \href{https://arxiv.org/abs/hep-ph/0503213}{arXiv:hep-ph/0503213}.

\bibitem[Erd\H{o}s \& R\'{e}nyi(1960)]{ErdosRenyi1960}
P.~Erd\H{o}s and A.~R\'{e}nyi.
\newblock On the evolution of random graphs.
\newblock \emph{Publications of the Mathematical Institute of the Hungarian
  Academy of Sciences}, 5:17--61, 1960.

\bibitem[Lovelock(1971)]{Lovelock1971}
D.~Lovelock.
\newblock The Einstein tensor and its generalizations.
\newblock \emph{Journal of Mathematical Physics}, 12(3):498--501, 1971.
\newblock \href{https://doi.org/10.1063/1.1665613}{doi:10.1063/1.1665613}.

\bibitem[Sandeman(2026a)]{Sandeman2026RAQM}
J.~F.~Sandeman.
\newblock Relational Actualism and Quantum Mechanics ({RAQM}).
\newblock \textit{Foundations of Physics} (Under review), 2026.
\newblock \href{https://doi.org/10.5281/zenodo.19174380}{doi.org/10.5281/zenodo.19174380}.

\bibitem[Sandeman(2026b)]{Sandeman2026RAGC}
J.~F.~Sandeman.
\newblock Relational Actualism: Gravity and Cosmology ({RAGC}).
\newblock Preprint, 2026.
\newblock \href{https://doi.org/10.5281/zenodo.19197999}{doi.org/10.5281/zenodo.19197999}.

\bibitem[Sandeman(2026c)]{Sandeman2026EB}
J.~F.~Sandeman.
\newblock The Engine of Becoming ({EB}).
\newblock Preprint, 2026.
\newblock \href{https://doi.org/10.5281/zenodo.19198120}{doi.org/10.5281/zenodo.19198120}.

\bibitem[Sandeman(2026d)]{Sandeman2026RAQI}
J.~F.~Sandeman.
\newblock Relational Actualism: Implications for Quantum Information ({RAQI}).
\newblock Preprint, 2026.
\newblock \href{https://doi.org/10.5281/zenodo.19198285}{doi.org/10.5281/zenodo.19198285}.

\bibitem[Sandeman(2026e)]{Sandeman2026RADM}
J.~F.~Sandeman.
\newblock Relational Actualism: Dark Matter as Actualization-Density Topology ({RADM}).
\newblock Preprint, 2026.
\newblock \href{https://doi.org/10.5281/zenodo.19198513}{doi.org/10.5281/zenodo.19198513}.

\bibitem[Glaser(2014)]{Glaser2014}
L.~Glaser.
\newblock A closed form expression for the causal set d'Alembertian.
\newblock \emph{Journal of Mathematical Physics}, 55(9):092301, 2014.
\newblock \href{https://doi.org/10.1063/1.4895604}{doi:10.1063/1.4895604}.
\newblock \href{https://arxiv.org/abs/1311.1701}{arXiv:1311.1701}.

\bibitem[Hawking(1975)]{Hawking1975}
S.~W.~Hawking.
\newblock Particle creation by black holes.
\newblock \emph{Communications in Mathematical Physics}, 43(3):199--220, 1975.
\newblock \href{https://doi.org/10.1007/BF02345020}{doi:10.1007/BF02345020}.

\bibitem[Page(1993)]{Page1993}
D.~N.~Page.
\newblock Information in black hole radiation.
\newblock \emph{Physical Review Letters}, 71(23):3743--3746, 1993.
\newblock \href{https://doi.org/10.1103/PhysRevLett.71.3743}{doi:10.1103/PhysRevLett.71.3743}.

\bibitem[Particle Data Group(2022)]{PDG2022}
R.~L.~Workman et~al.
\newblock Review of particle physics.
\newblock \emph{Progress of Theoretical and Experimental Physics},
  2022(8):083C01, 2022.
\newblock \href{https://doi.org/10.1093/ptep/ptac097}{doi:10.1093/ptep/ptac097}.

\bibitem[Peccei \& Quinn(1977)]{Peccei1977}
R.~D.~Peccei and H.~R.~Quinn.
\newblock CP conservation in the presence of pseudoparticles.
\newblock \emph{Physical Review Letters}, 38(25):1440--1443, 1977.
\newblock \href{https://doi.org/10.1103/PhysRevLett.38.1440}{doi:10.1103/PhysRevLett.38.1440}.

\bibitem[Sakharov(1967)]{Sakharov1967}
A.~D.~Sakharov.
\newblock Violation of CP invariance, C asymmetry, and baryon asymmetry of the universe.
\newblock \emph{JETP Letters}, 5:24--27, 1967.

\bibitem[Susskind(1993)]{Susskind1993}
L.~Susskind, L.~Thorlacius, and J.~Uglum.
\newblock The stretched horizon and black hole complementarity.
\newblock \emph{Physical Review D}, 48(8):3743--3761, 1993.
\newblock \href{https://doi.org/10.1103/PhysRevD.48.3743}{doi:10.1103/PhysRevD.48.3743}.

\bibitem[Susskind(2003)]{Susskind2003}
L.~Susskind.
\newblock The anthropic landscape of string theory.
\newblock \emph{arXiv preprint}, 2003.
\newblock \href{https://arxiv.org/abs/hep-th/0302219}{arXiv:hep-th/0302219}.

\end{thebibliography}

\end{document}
